\chapter{Value Numbering and Redundancy}
\label{ch:ch16-value-numbering-redundancy}

\section{Programs Repeat Themselves}
\label{sec:ch16-repeated-work}

Real programs compute the same value many times.
Address arithmetic reappears in every array reference inside a loop.
A common subexpression shows up once in the source and twice in the expanded IR.
Copy propagation and inlining create fresh duplicates that the programmer never wrote.
Each duplicate costs time, registers, and energy, and each one is an opportunity the compiler can take without changing what the program means.
The difficulty is recognizing duplicates reliably.
Textual identity is too weak, because \code{a + b} and \code{b + a} compute the same value with different text.
Textual identity is also too strong in the other direction, because two copies of \code{x + y} compute different values when \code{x} changes between them.
We need a notion of sameness that follows values through the program rather than comparing characters.
\term{Value numbering} gives every computed value a number, and two expressions share a number only when the compiler can prove they produce the same result.
A repeated expression with a known number is a candidate for \term{redundancy elimination}; replacement also needs a still-available value and compatible effects.
This chapter builds that signal from the ground up, starting inside a single basic block and then reaching across control flow.
The payoff is practical as well as conceptual.
Value numbering subsumes local common-subexpression elimination, copy propagation, and constant folding in one uniform mechanism.
It also prepares the ground for Chapter~\ref{ch:ch17-code-motion}, where expressions that are redundant on only some paths need code to move rather than merely disappear.

\section{Giving Equal Values Equal Numbers}
\label{sec:ch16-naming-values}

A \term{value number} names a runtime value rather than a storage location.
The variable \code{a} may hold different values at different points, so it carries different numbers at those points.
Two variables carry the same number at some point only when they are guaranteed to hold equal values there.
Constants receive numbers too, because folding \code{2 + 3} into \code{5} is the same phenomenon as recognizing that \code{a + b} appeared before.
The numbering starts by assigning fresh, distinct numbers to every variable and constant that has no known history.
Each new operation then forms a key from its operator and the value numbers of its operands.
When addition is commutative, the key sorts its operand numbers so that \code{a + b} and \code{b + a} form the same key.
A hash table maps each key seen so far to the number it was given.
A repeated key means a repeated value, and the compiler replaces the second computation with a copy from the first.
Replacements chain naturally, because a copied variable carries the number of its source.
The compiler asks one question at every expression, namely \enquote{have we computed this value before under these operand numbers}.
Everything else in the algorithm is bookkeeping that keeps that question honest.
The bookkeeping matters because a stale entry lies.
An assignment to \code{x} changes the value held in that variable, not values already computed from its former contents.
Expression keys over immutable value numbers can remain; a variable that has been overwritten must cease to represent its old number.
A store changes what a future load may return, not a scalar value already loaded into a temporary.
The local algorithm below handles pure, total scalar operations; memory numbering needs the additional state described in Chapter~\ref{ch:ch20-alias-analysis}.

\section{Numbering One Block at a Time}
\label{sec:ch16-local-numbering}

\term{Local value numbering} processes one basic block from top to bottom.
Straight-line code has no joins and no loops, so every earlier statement dominates every later one.
That property removes the dominance problem, but does not make the analysis complete for all algebraic equalities.
The algorithm keeps three maps while it scans.
One map sends each live expression key to its value number.
The second map sends each variable to the number of the value it currently holds.
The third records the current variable holders of each number, so replacement never reads an overwritten representative.
Listing~\ref{lst:ch16-block} shows a block that exercises every behavior of the algorithm.
It contains a genuine repetition, a repetition hidden by commutativity, a constant expression, a kill, and a repetition that survives the kill.

\begin{irlisting}[caption={A block with local redundancies, a kill, and a surviving repetition.},label={lst:ch16-block}]
; x and y hold unknown inputs on entry
t1 = x + y
t2 = y + x
t3 = 2 + 3
t4 = t1 * t3
x = 0
t5 = x + y
t6 = t1 * t3
\end{irlisting}

The first two lines compute the same sum in different orders, so the second line is redundant once commutativity is canonicalized.
The third line folds to the constant \code{5}.
The fourth line builds a product keyed by the numbers of \code{t1} and \code{5}.
The assignment to \code{x} removes \code{x} from the holders of its old number and gives it the number for zero.
The sum in line six therefore needs a fresh number, because its operand changed.
The product in line seven still matches, because neither \code{t1} nor \code{5} changed.
A purely textual comparison would miss the second line and would wrongly flag the sixth.
Value numbers get both cases right for the price of one hash lookup per expression.
Algorithm~\ref{alg:ch16-lvn} states the procedure precisely.

\begin{algorithm}
  \caption{Local value numbering with current value holders}\label{alg:ch16-lvn}
  \begin{pseudocode}
    \Require block $B$ of pure total scalar assignments with defined inputs
    \Ensure reuse only equal values with a current holder or literal representative
    \State $keynum \gets$ empty map from expression key to value number
    \State $varnum \gets$ empty map from variable to value number
    \State $holders \gets$ empty map from number to set of current variables
    \For{each instruction $I$ in $B$ in order}
      \State record destination $v$; give unknown input variables fresh numbers and register their holders
      \State read operand numbers before changing $varnum[v]$
      \State fold literal-only operations; intern typed literals by value
      \State copies use their source number; other operations use a canonical typed key
      \State $n \gets$ source number for a copy; otherwise look up or intern the value's key in $keynum$
      \If{$n$ has a literal representative or a nonempty holder set}
        \State replace $I$ with that literal or a copy from a current holder of $n$
      \EndIf
      \State remove $v$ from its former holder set, if any
      \State $varnum[v] \gets n$; add $v$ to $holders[n]$
    \EndFor
  \end{pseudocode}
\end{algorithm}

The algorithm assumes bounded-arity, pure total operations and no hidden definitions.
Keys include result type and semantic flags; commutativity is used only when valid for that operation.
Its invariant is that every variable in $holders[n]$ currently contains the value denoted by $n$.
Operand numbers are read before the destination mapping changes, so \code{x = x + y} is handled correctly.
If a key is known but has no surviving holder, keep the computation and register its destination as a new holder of the old number.
For example, after \code{t = a + b; t = 0}, a later \code{u = a + b} cannot become \code{u = t}.
Across an unmodeled effect, conservatively end the numbering region and start fresh; extending the region requires explicit memory and exception semantics.

\section{Worked Example: Numbering the Block}
\label{sec:ch16-worked-numbering}

We now run Algorithm~\ref{alg:ch16-lvn} on Listing~\ref{lst:ch16-block} by hand.
Inputs \code{x} and \code{y} receive the fresh numbers $1$ and $2$.
The key for line one is $plus(1, 2)$, which is new, so it receives number $3$ and \code{t1} maps to $3$.
Line two builds $plus(2, 1)$, and canonicalization sorts the operands into $plus(1, 2)$.
That key is present with number $3$, so line two becomes the copy \code{t2 = t1} and \code{t2} maps to $3$.
Line three folds the literal expression $2+3$ into the constant \code{5}, which receives the fresh number $4$, and \code{t3} maps to $4$.
Line four builds $mul(3, 4)$, which is new, so it receives number $5$ and \code{t4} maps to $5$.
Line five assigns the constant \code{0}, which receives number $6$, and \code{x} now maps to $6$.
The holder update removes \code{x} from number $1$; the key $plus(1,2)$ remains valid and its result is still held by \code{t1} and \code{t2}.
Line six builds $plus(6, 2)$, which is absent, so it receives number $7$ and \code{t5} maps to $7$.
Line seven builds $mul(3,4)$ and finds number $5$ still held by \code{t4}, so it becomes the copy \code{t6 = t4}.
Table~\ref{tab:ch16-trace} records the full trace for checking.

\begin{table}
  \centering
  \caption{Value numbers assigned while scanning Listing~\ref{lst:ch16-block} from top to bottom.}
  \label{tab:ch16-trace}
  \begin{tabular}{llll}
    \toprule
    Line & Instruction & Key & Action \\
    \midrule
    1 & \code{t1 = x + y} & $plus(1,2)$ & new number 3 \\
    2 & \code{t2 = y + x} & $plus(1,2)$ & reuse 3, copy \code{t1} \\
    3 & \code{t3 = 2 + 3} & const 5 & new number 4 \\
    4 & \code{t4 = t1 * t3} & $mul(3,4)$ & new number 5 \\
    5 & \code{x = 0} & const 0 & new number 6, update holder \\
    6 & \code{t5 = x + y} & $plus(6,2)$ & new number 7 \\
    7 & \code{t6 = t1 * t3} & $mul(3,4)$ & reuse 5, copy \code{t4} \\
    \bottomrule
  \end{tabular}
\end{table}

Two replacements survive, one folding happens, and one near miss is correctly rejected.
The rewritten block computes three fewer operations than the original.
Updating the mapping of \code{x} in line five is essential.
The new number $6$ makes the next addition's key differ from $plus(1,2)$ even if the old expression key remains in the table.
Deleting that old key is conservative rather than necessary in this value-numbered example; reusing an overwritten representative variable would be the real hazard.
When in doubt, kill more and let a later pass clean up, because a wrong reuse miscompiles while a missed reuse merely runs slower.

\section{Reaching Across Basic Blocks}
\label{sec:ch16-global-numbering}

Local numbering stops at block boundaries, but redundancy does not.
Figure~\ref{fig:ch16-join} shows the classic cross-block shape.
The entry block computes \code{a + b}, both branches compute it again, and the join block computes it a third time.
The computation in the join block is redundant only when every path into the join carries the same value.
Here both branches recompute rather than inherit, yet each branch computation is itself redundant with the entry.
Local numbering alone cannot see the entry computation from another block; a dominator-scoped or global pass can replace all three later additions with the entry value.
\term{Global value numbering} reasons about all paths at once instead of rescanning each block in isolation.
The standard formulation runs over SSA form, where each variable has a single definition and each join carries explicit \code{phi} functions.
Matching operations and operand numbers can establish equal SSA values; phi comparison must also preserve the correspondence between inputs and incoming edges.
Equality does not imply availability: a replacement definition must dominate the use, or the pass must construct a suitable phi from edge-available values.
Congruence propagates through the SSA graph the way it propagates down a block, except that joins need a fixed point.
% TODO(cite): primary reference for SSA-based global value numbering and its congruence closure.
An expression is \term{available} at a point when every path from the entry computes it without redefining its operands afterward.
Availability is a forward dataflow fact, and the join-block computation above is redundant exactly when its expression is available there.
Table~\ref{tab:ch16-kinds} classifies the three shapes of redundancy and the remedy each one needs.

\begin{figure}
  \centering
  \begin{tikzpicture}[node distance=14mm and 18mm]
    \node[basic block] (ch16b0) {t0 = a + b};
    \node[basic block,below left=of ch16b0] (ch16b1) {t1 = a + b};
    \node[basic block,below right=of ch16b0] (ch16b2) {t2 = a + b};
    \node[basic block,below right=of ch16b1] (ch16b3) {t3 = a + b};
    \draw[cfg edge] (ch16b0) -- (ch16b1);
    \draw[cfg edge] (ch16b0) -- (ch16b2);
    \draw[cfg edge] (ch16b1) -- (ch16b3);
    \draw[cfg edge] (ch16b2) -- (ch16b3);
  \end{tikzpicture}
  \caption{The join block recomputes a sum available on every path, so it is fully redundant once the branches reuse the entry value.}
  \label{fig:ch16-join}
\end{figure}

\begin{table}
  \centering
  \caption{Three shapes of redundancy differ in where the earlier computation sits and what fix each one needs.}
  \label{tab:ch16-kinds}
  \begin{tabular}{lll}
    \toprule
    Shape & Where the duplicate sits & Remedy \\
    \midrule
    Local & Same basic block & Copy from the first computation \\
    Global & Every path to the use & Dataflow availability plus reuse \\
    Partial & Some paths to the use & Insert on missing paths, then reuse \\
    \bottomrule
  \end{tabular}
\end{table}

The third row is deliberately unresolved here.
Partial redundancy needs code to move, not just code to vanish, and Chapter~\ref{ch:ch17-code-motion} takes up that problem.
Global numbering also interacts with the kills discussed earlier.
A call or a possibly aliasing store inside one branch ends availability for expressions it might redefine.
Precision therefore depends on the alias answers of Chapter~\ref{ch:ch20-alias-analysis}, and production compilers rerun numbering after inlining precisely because new paths create new redundancies.

\section{Implementation Notes}
\label{sec:ch16-implementation}

For $n$ bounded-arity instructions, the scalar algorithm takes expected $O(n)$ time and $O(n)$ space with hash tables and constant-time holder insertion, removal, and representative selection.
Store a direct holder-list link for each variable so overwriting a variable does not require searching a set.
No expression-table scan is needed on a scalar assignment.
This bound assumes bounded-size keys and constant-time arithmetic; arbitrary-precision folding needs its own arithmetic cost.
Commutative canonicalization should sort operand numbers once at key construction.
Constant folding belongs inside key construction too, so that \code{2 + 3} and \code{5} meet in the table rather than in a separate pass.
Associative and distributive laws are tempting but dangerous.
Reassociation changes rounding for floating-point values and changes trapping behavior for integers near overflow, so keep the key builder to identities the language guarantees.
Production compilers run value numbering on SSA with hash consing, which shares one table entry per distinct expression across the whole function.
% TODO(cite): modern production GVN designs and their predicate handling.
They may use predicates from branch conditions, so a path known to have \code{x} nonzero admits simplifications unavailable elsewhere.
Memory expressions need versioning against stores, which turns loads into keyed expressions over an abstract memory state.
Global complexity depends on the chosen congruence solver, loop processing, and memory model; the local linear bound does not establish a bound for those extensions.
Start with the local algorithm, add SSA-based global numbering, and only then reach for predicates and memory versioning.

\section{From congruence classes to retained alternatives}
\label{sec:ch16-egraph-link}

Value numbering recognizes equal computations already represented in the program and asks whether an available value can replace a redundant one.
Equality saturation broadens the search: it adds soundly equivalent expressions, retains several alternatives, and chooses a representative later.
For $(x+0)\cdot1$ in pure byte arithmetic, identity rules let one class represent both the original tree and the single value $x$.
This class is not a location from which generated code can simply load; extraction must produce a finite executable expression and later placement must satisfy dominance.

Chapter~\ref{ch:equality-saturation} implements this distinction with class union, congruence rebuilding, and positive-cost extraction.
The local value-numbering key discipline still matters: types, arithmetic semantics, and effects determine which equalities may enter the graph.
A cheaper representative selected by a cost model is useful only after those semantic conditions have been established.

\section{Exercises}
\label{sec:ch16-exercises}

\begin{enumerate}
  \item \label{ex:ch16-trace-kill}
  Number the following block by hand with Algorithm~\ref{alg:ch16-lvn}.
  Give the value number of every variable after each line, state which lines become copies, and show the final rewritten block.
  \code{p = a + b}, \code{q = p + c}, \code{a = b}, \code{r = a + b}, \code{s = p + c}.
  Explain why overwriting \code{a} does not invalidate the value already held in \code{p}, and identify the surviving holder used for the final replacement.
  \item \label{ex:ch16-commutative-miss}
  Give a three-line block where a purely textual common-subexpression pass finds nothing but value numbering with canonicalization removes one computation.
  Run both the textual pass and Algorithm~\ref{alg:ch16-lvn} on your block and show the two different outputs.
  Then insert an assignment that changes a shared operand before the repeated expression, rerun the algorithm, and explain why its key changes.
  \item \label{ex:ch16-global-availability}
  Draw a CFG with an entry block, two branch blocks, and a join block where the entry computes \code{m * n}, only the left branch recomputes it, and the join computes it again.
  Compute the available expressions at the join by hand and decide whether the join computation is fully redundant.
  Explain why no insertion is needed while the entry computation remains.
  Then remove the entry computation, recompute availability, and identify the missing-path insertion needed for partial redundancy elimination, assuming pure total multiplication.
\end{enumerate}

\section{Further reading}
\label{sec:ch16-further-reading}

Morgan's \emph{Building an Optimizing Compiler} treats redundancy elimination as a core optimization with attention to its dataflow foundations~\cite{morgan-1998-optimizing}.
Read it for how availability and related facts connect local rewriting to whole-function reasoning.
Appel's \emph{Modern Compiler Implementation in ML} places these ideas inside a complete compiler where ordering among passes matters~\cite{appel-1998-modern}.
Use it to see where numbering fits relative to SSA construction and later cleanup passes.
Safonov's \emph{Trustworthy Compilers} discusses what must hold for an optimization to preserve meaning~\cite{safonov-2010-trustworthy}.
Consult it before extending the key builder with algebraic laws that the source language does not guarantee.
